\section{A Model for a Sequence of Independent Random Variables}

In this section we show how to construct a countably infinite sequence of independent random variables by drawing a random number uniformly in the unit interval $[0,1]$. The uniform distribution is a fundamental building block in probability theory, as it contains enough randomness to generate an entire sequence of independent random variables. With this probability model, we prove the Kolmogorov zero--one law.

For any real number $u$ between 0 and 1, we let $B_i(u)$ be the $i$-th bit in the binary expansion of $u$, i.e.,
\[
u = \sum_{i=1}^{\infty} B_i(u) 2^{-i}.
\]

When a number has more than one binary expansion, we choose the one with infinitely many trailing zeros. For any real number $u$, we obtain an infinite binary sequence $(B_i(u))_{i\ge 1}$. Note that the sequences with infinitely many trailing ones are not represented in this probability model. However, these missing sequences form a set of probability $0$ and are therefore negligible.

The first bit $B_1(u)$ is equal to $0$ if $u \in A_1 = [0,1/2)$ and $1$ if $u \in [1/2,1)$. The second bit $B_2(u)$ equals $0$ if $u \in A_2 = [0,1/4)\cup [1/2,3/4)$ and $1$ if $u$ is in the complement of $A_2$ in $[0,1)$. In general, the first $n$ bits of $u$ are $x_1, x_2, \dots, x_n$, where $x_i$ are binary digits in $\{0,1\}$, when
\[
\sum_{i=1}^{n} x_i 2^{-i} \le u < \frac{1}{2^n} + \sum_{i=1}^{n} x_i 2^{-i}.
\tag{5.5}
\]

The mapping from a real number $u$ to the associated sequence $(B_i(u))_{i\ge 1}$ described above is a deterministic procedure. Specifically, $B_i(u)$ is the $i$-th bit in the binary expansion of $u$. We now construct a Lebesgue--Stieltjes measure on $[0,1]$ using the Stieltjes measure function $F(x)=x$, which is the identity function, for $0\le x\le 1$. This measure, denoted by $P$, defines a probability space $([0,1], \mathcal{B}([0,1]), P)$.

The $i$-th bit $B_i(u)$ becomes a random variable under the probability measure $P$. For $x_1,x_2,\dots,x_n \in \{0,1\}$, the probability that the bit sequence $B_i=x_i$ for $i=1,2,\dots,n$ is
\[
P(B_1 = x_1, B_2 = x_2, \dots, B_n = x_n) = \frac{1}{2^n}.
\tag{5.6}
\]

By (5.5), this process generates a sequence of independent Bernoulli random variables, each with probability $1/2$ of being 0 or 1, and the joint distribution of these random variables is uniform on the set of all binary sequences with length $n$.

\medskip

We divide the bits into two streams. The first stream consists of the bits $B_i(u)$ with even indices $i$ and the second with odd indices $i$,
\[
S_e(u) \triangleq (B_2(u), B_4(u), B_6(u), \dots), \quad S_o(u) \triangleq (B_1(u), B_3(u), B_5(u), \dots).
\]

We can define two new random variables $X_e(u)$ and $X_o(u)$ by
\[
X_e(u) \triangleq \sum_{i=1}^{\infty} B_{2i}(u) 2^{-i}, \quad X_o(u) \triangleq \sum_{i=1}^{\infty} B_{2i-1}(u) 2^{-i}.
\]

The two random variables $X_e$ and $X_o$ are independent and uniformly distributed on $[0,1]$. They are independent because the bits with even indices and the bits with odd indices are independent. We may call this process ``splitting.''

We let $U_1$ to be the random variable $X_e$. We then split the bits stream $S_o$ into odd and even parts and use one part to define a uniform random variable $U_2$. This process can be repeated recursively to produce an infinite sequence of independent uniform random variables.

We have thus constructed an infinite sequence $(U_i(u))_{i\ge 1}$ of i.i.d. uniform random variables from a uniformly distributed random variable. By using the inverse transform method (see Problem 5.7), we can transform this sequence to an infinite sequence of independent random variables with arbitrary distribution. For example, if we want a sequence of independent but biased coin tosses, we can proceed as follows. For $i=1,2,\dots$, let $Y_i$ be a binary random variable defined by $Y_i(u) = \mathbf{1}_{\{U_i(u) < p_i\}}$, where $p_1,p_2,p_3,\dots$ are given constants. The sequence $(Y_i(u))_{i\ge 1}$ represents a sequence of biased but independent coin tosses.

It is important to note that once a sample point $u$ between 0 and 1 is selected, the whole sequence $(Y_i(u))_{i\ge 1}$ is determined. This construction procedure provides an explicit probability model for an infinite sequence of independent random variables.

Kolmogorov's zero--one law states that the conclusion in the Borel zero--one law is a common phenomenon. We will use the notation $\sigma(X_1,X_2,X_3,\dots)$ to denote the smallest $\sigma$-algebra under which the random variables $X_i$'s are all measurable.

\begin{theorem}[Kolmogorov's Zero--One Law]
Let $(X_n)_{n=1}^\infty$ be a sequence of random variables, and define the tail $\sigma$-algebra by $\mathcal{T} \triangleq \bigcap_{n=1}^{\infty} \sigma(\{X_i:i\ge n\})$. If the random variables $X_n$, for $n\ge 1$, are mutually independent, then any event in the tail $\sigma$-algebra has probability equal to either $0$ or $1$.
\end{theorem}

\begin{proof}
The idea of proof is to show that the tail $\sigma$-algebra $\mathcal{T}$ is independent of itself. Let $\mathcal{G}_n$ denote the $\sigma$-algebra generated by the first $n$ random variables $X_1,\dots,X_n$ in the sequence. The $\sigma$-algebra $\mathcal{G}_n$ is independent of the tail $\sigma$-algebra $\mathcal{T}$, because any event in $\mathcal{T}$ is in $\sigma(\{X_i:i>n\})$, which is independent with $\mathcal{G}_n$.
\end{proof}
